\documentclass{article}
\usepackage[utf8]{inputenc}
\usepackage{tikz,readarray}
\usepackage[a4paper]{geometry}
\geometry{a4paper,landscape,left=30mm,top=20mm,right=10mm,bottom=10mm}
\begin{document}
\pagestyle{empty}
\readarraysepchar{\\}%
\readdef{2025-4117-5-3.csv}\mydata%
\ignoreemptyitems
\setsepchar{\\/,}
\readlist\mylist{\mydata}%




\begin{tikzpicture}[remember picture, overlay,shift={(-2.5,-18)}]
\foreach \c in {0,...,0}
\foreachitem\x\in\mylist{%
  \draw (0,17.7-\xcnt*0.5+1.5-\c*6.5) node [right]{\x};
};
%\draw [dashed](2,19) --(27,19); 
%\draw [dashed](2,19.5) --(27,19.5); 
\foreach \z in {0,...,47}
     \draw (3+\z*0.5,18)--(3.5+\z*0.5,18.5);

\draw (0,18.25) node [right] {Apellido y Nombre};

\draw [step=0.5,thin,xstep=5.5] (0,0) grid (2.5,18.5);
\draw [step=0.5,thin] (2.5,0) grid (27,18.5);
\draw [thin] (0,0) rectangle (27,18.5);
\draw [ultra thick](0,18)--(27,18);

\end{tikzpicture}

% Paginas con notas manuscritas
\newpage
\begin{tikzpicture}[remember picture, overlay,shift={(-2.5,-18)}]
\draw (3,19) node []{\huge Notas};
\draw(0,0) rectangle (27.5,18);
\foreach \y in {-0,...,36}
        \draw (0,\y/2) --(27.5,\y/2);

\end{tikzpicture}

    
    
\end{document}